% Part: first-order-logic
% Chapter: natural-deduction
% Section: proving-things-quant

\documentclass[../../../include/open-logic-section]{subfiles}

\begin{document}

\olfileid{fol}{ntd}{prq}

\olsection{\usetoken{P}{derivation} kanthi Kuantor}

\begin{ex}
Nalika nangani kuantor, kita kudu mesthekake supaya ora nerak syarat
eigenvariabel, lan kadhangkala iki mbutuhake kita ngowahi urutan
panrapan inferensi tartamtu. Umume, luwih gampang yen kita nyoba
nangani luwih dhisik aturan kang kena syarat eigenvariabel
(aturan-aturan iki bakal ana ing sisih ngisor bukti kang wis rampung).

Ayo ndeleng carane menehi !!a{derivation} saka !!{formula}
$\lexists[x][\lnot !A(x)] \lif \lnot \lforall[x][!A(x)]$.
Kaya lumrahe, kita miwiti kanthi nulis
\begin{prooftree}
\AxiomC{}
\UnaryInfC{$\lexists[x][\lnot !A(x)]\lif \lnot \lforall[x][!A(x)]$}
\end{prooftree}
Kita miwiti kanthi nulis apa kang dibutuhake kanggo mbenerake langkah
pungkasan nganggo aturan \Intro{\lif}.
\begin{prooftree}
\AxiomC{$\Discharge{\lexists[x][\lnot !A(x)]}{1}$}
\DeduceC{$\lnot \lforall[x][!A(x)]$}
\DischargeRule{\Intro{\lif}}{1}
\UnaryInfC{$\lexists[x][\lnot !A(x)]\lif \lnot \lforall[x][!A(x)]$}
\end{prooftree}
Amarga ora ana aturan cetha kang bisa ditrapake marang $\lnot
\lforall[x][!A(x)]$, kita bakal nata !!{derivation} supaya bisa
nggunakake aturan \Elim{\lexists}. Ing kene kita kudu nggatekake
syarat eigenvariabel lan milih konstanta kang ora muncul ing
$\lexists[x][\lnot !A(x)]$ utawa ing asumsi apa wae kang dadi
gantungane. Cathetan koreksi sumber OLPL-012: formula Inggris ing
ukara iki ngilangake negasi kang ana ing asumsi eksistensial lan saben
wit ing tuladha; wujud kang dibalekake nduweni himpunan konstanta kang
padha, mula syarat kesegaran ora owah. (Nanging, amarga ora ana
!!{constant}s kang muncul, pilihan apa wae bisa digunakake.)
\begin{prooftree}
\AxiomC{$\Discharge{\lexists[x][\lnot !A(x)]}{1}$}
\AxiomC{$\Discharge{\lnot !A(a)}{2}$}
\DeduceC{$\lnot \lforall[x][!A(x)]$}
\DischargeRule{\Elim{\lexists}}{2}
\BinaryInfC{$\lnot \lforall[x][!A(x)]$}
\DischargeRule{\Intro{\lif}}{1}
\UnaryInfC{$\lexists[x][\lnot !A(x)] \lif \lnot \lforall[x][!A(x)]$}
\end{prooftree}
Kanggo ngasilake $\lnot \lforall[x][!A(x)]$, kita bakal nyoba
nggunakake aturan \Intro{\lnot}: iki mbutuhake kita ngasilake
kontradiksi, bisa uga kanthi nggunakake $\lforall[x][!A(x)]$ minangka
asumsi tambahan. Mesthi wae, kontradiksi iki bisa nyakup asumsi
$\lnot !A(a)$ kang bakal dibubarake dening inferensi
\Elim{\lexists}. Kita bisa nata bukti kaya mangkene:
\begin{prooftree}
\AxiomC{$\Discharge{\lexists[x][\lnot !A(x)]}{1}$}
\AxiomC{$\Discharge{\lnot !A(a)}{2}, \Discharge{\lforall[x][!A(x)]}{3}$}
\DeduceC{$\lfalse$}
\DischargeRule{\Intro{\lnot}}{3}
\UnaryInfC{$\lnot \lforall[x][!A(x)]$}
\DischargeRule{\Elim{\lexists}}{2}
\BinaryInfC{$\lnot \lforall[x][!A(x)]$}
\DischargeRule{\Intro{\lif}}{1}
\UnaryInfC{$\lexists[x][\lnot !A(x)]\lif \lnot \lforall[x][!A(x)]$}
\end{prooftree}
Katone kita wis meh ngasilake kontradiksi. Aturan kang paling gampang
ditrapake yaiku \Elim{\lforall}, kang ora nduweni syarat
eigenvariabel. Amarga kita bisa nggunakake sembarang term kanggo
ngganti $x$ kang dikuantifikasi universal, paling trep yen kita terus
nggunakake~$a$ supaya bisa nggayuh kontradiksi.
\begin{prooftree}
\AxiomC{$\Discharge{\lexists[x][\lnot !A(x)]}{1}$}
\AxiomC{$\Discharge{\lnot !A(a)}{2}$}
\AxiomC{$\Discharge{\lforall[x][!A(x)]}{3}$}
\RightLabel{\Elim{\lforall}}
\UnaryInfC{$!A(a)$}
\RightLabel{\Elim{\lnot}}
\BinaryInfC{$\lfalse$}
\DischargeRule{\Intro{\lnot}}{3}
\UnaryInfC{$\lnot \lforall[x][!A(x)]$}
\DischargeRule{\Elim{\lexists}}{2}
\BinaryInfC{$\lnot \lforall[x][!A(x)]$}
\DischargeRule{\Intro{\lif}}{1}
\UnaryInfC{$\lexists[x][\lnot !A(x)]\lif \lnot \lforall[x][!A(x)]$}
\end{prooftree}

Ing tahap iki, mligine nalika nangani kuantor, kita kudu mriksa maneh
supaya syarat eigenvariabel ora dilanggar. Amarga siji-sijine aturan
kang kita trapake lan kena syarat eigenvariabel yaiku
\Elim{\exists}, lan eigenvariabel~$a$ ora muncul ing asumsi apa wae
kang dadi gantungane, iki !!{derivation} kang bener.
\end{ex}

\begin{ex}
Kadhangkala kita bisa ngasilake !!a{formula} saka !!{formula}s liyane.
Ing kasus iki bisa ana asumsi kang durung dibubarake. Mula kita kudu
nyathet asumsi-asumsi kasebut uga ancas pungkasane.

Ayo ndeleng carane menehi !!a{derivation} saka !!{formula}
$\lexists[x][!C(x,b)]$ adhedhasar asumsi $\lexists[x][(!A(x)
\land !B(x))]$ lan $\lforall[x][(!B(x) \lif !C(x,b))]$.
Kaya lumrahe, kita nulis kesimpulan ing sisih ngisor.
\begin{prooftree}
\AxiomC{}
\UnaryInfC{$\lexists[x][!C(x,b)]$}
\end{prooftree}

Kita duwe rong premis kang bisa digunakake. Kanggo nggunakake premis
kapisan, yaiku nyoba nemokake !!a{derivation} saka
$\lexists[x][!C(x, b)]$ adhedhasar $\lexists[x][(!A(x)
  \land !B(x))]$, kita bakal nggunakake aturan \Elim{\lexists}.
Amarga aturan kasebut nduweni syarat eigenvariabel, kita bakal
ngetrapake aturan iku luwih dhisik. Asile mangkene:
\begin{prooftree}
\AxiomC{$\lexists[x][(!A(x) \land !B(x))]$}
\AxiomC{$\Discharge{!A(a) \land !B(a)}{1}$}
\DeduceC{$\lexists[x][!C(x,b)]$}
\DischargeRule{\Elim{\lexists}}{1}
\BinaryInfC{$\lexists[x][!C(x,b)]$}
\end{prooftree}
Rong asumsi kang kita gunakake padha-padha ngemot~$!B$. Ing tahap iki
bisa migunani yen \Elim{\land} ditrapake kanggo misahake~$!B(a)$.
\begin{prooftree}
\AxiomC{$\lexists[x][(!A(x) \land !B(x)])$}
\AxiomC{$\Discharge{!A(a)
\land !B(a)}{1}$}
\RightLabel{\Elim{\land}}
\UnaryInfC{$!B(a)$}
\DeduceC{$\lexists[x][!C(x,b)]$}
\DischargeRule{\Elim{\lexists}}{1}
\BinaryInfC{$\lexists[x][!C(x,b)]$}
\end{prooftree}

Asumsi kapindho kang bisa digunakake yaiku~$\lforall[x][(!B(x) \lif
  !C(x,b))]$. Amarga ora ana syarat eigenvariabel, kita bisa
ngganti $x$ nganggo !!{constant}~$a$ liwat \Elim{\lforall} kanggo
ngasilake $!B(a) \lif !C(a, b)$. Saiki kita duwe $!B(a) \lif !C(a,b)$
lan $!B(a)$. Langkah sabanjure yaiku panrapan langsung saka aturan
\Elim{\lif}.
\begin{prooftree}
\AxiomC{$\lexists[x][(!A(x) \land !B(x))]$}
\AxiomC{$\lforall[x][(!B(x) \lif !C(x,b))]$}
\RightLabel{\Elim{\lforall}}
\UnaryInfC{$!B(a) \lif !C(a,b)$}
\AxiomC{$\Discharge{!A(a)
\land !B(a)}{1}$}
\RightLabel{\Elim{\land}}
\UnaryInfC{$!B(a)$}
\RightLabel{\Elim{\lif}}
\BinaryInfC{$!C(a,b)$}
\DeduceC{$\lexists[x][!C(x,b)]$}
\DischargeRule{\Elim{\lexists}}{1}
\insertBetweenHyps{\hspace{-5em}}
\BinaryInfC{$\lexists[x][!C(x,b)]$}
\end{prooftree}
Wis cedhak banget!{} Siji panrapan \Intro{\lexists} maneh lan kita
wis nggayuh ancas.
\begin{prooftree}
\AxiomC{$\lexists[x][(!A(x) \land !B(x))]$}
\AxiomC{$\lforall[x][(!B(x) \lif !C(x,b))]$}
\RightLabel{\Elim{\lforall}}
\UnaryInfC{$!B(a) \lif !C(a,b)$}
\AxiomC{$\Discharge{!A(a)
\land !B(a)}{1}$}
\RightLabel{\Elim{\land}}
\UnaryInfC{$!B(a)$}
\RightLabel{\Elim{\lif}}
\BinaryInfC{$!C(a,b)$}
\RightLabel{\Intro{\lexists}}
\UnaryInfC{$\lexists[x][!C(x,b)]$}
\DischargeRule{\Elim{\lexists}}{1}
\insertBetweenHyps{\hspace{-5em}}
\BinaryInfC{$\lexists[x][!C(x,b)]$}
\end{prooftree}
Amarga ing saben langkah kita wis mesthekake supaya syarat
eigenvariabel ora dilanggar, kita bisa yakin yen iki !!{derivation}
kang bener.
\end{ex}

\begin{ex}
Wènèhana !!a{derivation} saka !!{formula}
$\lnot\lforall[x][!A(x)]$ adhedhasar asumsi $\lforall[x][!A(x)]
\lif \lexists[y][!B(y)]$ lan $\lnot\lexists[y][!B(y)]$.
Kaya lumrahe, kita nulis !!{formula} sasaran ing sisih ngisor.
\begin{prooftree}
\AxiomC{}
\UnaryInfC{$\lnot\lforall[x][!A(x)]$}
\end{prooftree}
Larik pungkasan !!{derivation} iku negasi, mula ayo nyoba nggunakake
\Intro{\lnot}. Iki mbutuhake kita nemokake cara kanggo !!{derive}
kontradiksi.
\begin{prooftree}
\AxiomC{$\Discharge{\lforall[x][!A(x)]}{1}$}
\DeduceC{$\lfalse$}
\DischargeRule{\Intro{\lnot}}{1}
\UnaryInfC{$\lnot\lforall[x][!A(x)]$}
\end{prooftree}
Nganti kene kabeh lumaku becik. Kita bisa nggunakake \Elim{\lforall},
nanging durung cetha apa aturan iku bakal mbantu nggayuh ancas.
Mula, ayo nggunakake salah siji asumsi kita.
$\lforall[x][!A(x)] \lif \lexists[y][!B(y)]$ bebarengan karo
$\lforall[x][!A(x)]$ ngidini kita nggunakake aturan \Elim{\lif}.
\begin{prooftree}
\AxiomC{$\lforall[x][!A(x)] \lif \lexists[y][!B(y)]$}
\AxiomC{$\Discharge{\lforall[x][!A(x)]}{1}$}
\RightLabel{\Elim{\lif}}
\BinaryInfC{$\lexists[y][!B(y)]$}
\DeduceC{$\lfalse$}
\DischargeRule{\Intro{\lnot}}{1}
\UnaryInfC{$\lnot\lforall[x][!A(x)]$}
\end{prooftree}
Saiki isih ana siji asumsi pungkasan kang bisa digunakake, lan
katone asumsi iki bakal mbantu kita nggayuh kontradiksi liwat
\Elim{\lnot}.
\begin{prooftree}
\AxiomC{$\lnot\lexists[y][!B(y)]$}
\AxiomC{$\lforall[x][!A(x)] \lif \lexists[y][!B(y)]$}
\AxiomC{$\Discharge{\lforall[x][!A(x)]}{1}$}
\RightLabel{\Elim{\lif}}
\BinaryInfC{$\lexists[y][!B(y)]$}
\RightLabel{\Elim{\lnot}}
\BinaryInfC{$\lfalse$}
\DischargeRule{\Intro{\lnot}}{1}
\UnaryInfC{$\lnot\lforall[x][!A(x)]$}
\end{prooftree}
\end{ex}

\begin{prob}
Wènèhana !!{derivation}s kang nuduhake prakara-prakara iki:
\begin{enumerate}
\item $\Proves (\lforall[x][!A(x)] \land \lforall[y][!B(y)]) \lif
\lforall[z][(!A(z) \land !B(z))]$.
\item $\Proves (\lexists[x][!A(x)] \lor \lexists[y][!B(y)]) \lif
\lexists[z][(!A(z) \lor !B(z))]$.
\item $\lforall[x][(!A(x) \lif !B)] \Proves \lexists[y][!A(y)] \lif !B$.
\item $\lforall[x][\lnot !A(x)] \Proves \lnot\lexists[x][!A(x)]$.
\item $\Proves \lnot\lexists[x][!A(x)] \lif \lforall[x][\lnot !A(x)]$.
\item $\Proves \lnot\lexists[x][\lforall[y][((!A(x,y) \lif \lnot
!A(y,y)) \land (\lnot !A(y,y) \lif !A(x,y)))]]$.
\end{enumerate}
\end{prob}

\begin{prob}
Wènèhana !!{derivation}s kang nuduhake prakara-prakara iki:
\begin{enumerate}
\item $\Proves \lnot\lforall[x][!A(x)] \lif \lexists[x][\lnot!A(x)]$.
\item $(\lforall[x][!A(x)] \lif !B) \Proves \lexists[y][(!A(y) \lif !B)]$.
\item $\Proves \lexists[x][(!A(x) \lif \lforall[y][!A(y)])]$.
\end{enumerate}
(Kabeh iki mbutuhake aturan~$\FalseCl$.)
\end{prob}

\end{document}
